\chapter{Análise e Discussão}\label{cap:discussao}

Neste capítulo, interpretaram-se as diferenças entre o coeficiente de Pearson $r$
e o escore de coerência $S_{xy}^{(0{,}05)}$, a sensibilidade dos resultados e
seu alcance frente à literatura. A discussão tomou como referência as 26 aquisições
e todos os 325 pares do \autoref{cap:resultados}.

% =========================================================
\section{Interpretação das duas medidas}\label{sec:disc-interpretacao}
% =========================================================

O contraste descritivo $\Delta_h=0{,}0659$ do coeficiente de Pearson $r$, com peso igual
por par, indicou maior semelhança espectral média dentro dos padrões. O contraste
médio $\Delta_h=0{,}1453$, com peso igual por grupo, preservou essa direção.
Esse comportamento constituiu um indício descritivo promissor para a investigação
de assinaturas espectrais, principalmente pela proximidade dos descritores de
\textit{plug} e de parte dos ensaios anulares. A ampla sobreposição da
\autoref{fig:res-distribuicoes}, entretanto, impediu interpretar o contraste
como separação automática dos padrões ou como significância estatística demonstrada.

O coeficiente de Pearson $r$ comparou a forma dos resíduos espectrais, conforme
a \autoref{eq:met-pearson}. Em termos práticos, dois descritores receberam
coeficiente de Pearson $r$ elevado quando suas elevações e reduções relativas
ocorreram em comprimentos de onda semelhantes, depois da transformação logarítmica
e da remoção da tendência. Assim, a boa concordância dos descritores não exigiu
a coincidência das estruturas em uma mesma posição do panorama. Também não
preservou a amplitude original nem a inclinação espectral removida.

O escore de coerência $S_{xy}^{(0{,}05)}$ resumiu outra propriedade: os máximos
da coerência wavelet encontrados nos atrasos de cada escala, ponderados pela potência
dos dois sinais. Além disso, utilizou prefixos cuja extensão dependeu do parceiro
da comparação. Por isso, as duas medidas não responderam à mesma pergunta.
O menor escore de coerência intrapadrão não invalidou a semelhança dos resíduos
espectrais; indicou que esse resumo dos campos não favoreceu a discriminação
global pretendida.

A \autoref{fig:disc-inter} tornou essa distinção concreta. O par m04q01, de
\textit{slug}, e m07q20, de \textit{churn}, apresentou escore de coerência
$S_{xy}^{(0{,}05)}=0{,}5554$ e máximo da coerência wavelet de $0{,}9950$,
apesar dos rótulos diferentes. O coeficiente de Pearson $r$ desse par foi
$0{,}7103$. Uma região clara no mapa não bastou, portanto, para atribuir
igualdade de padrão aos dois ensaios. Esse exemplo ilustrou o alcance da medida;
a avaliação global permaneceu baseada em todos os pares.

\begin{figure}[p]
  \centering
  \includegraphics[width=\textwidth,height=0.86\textheight,keepaspectratio]{figuras/v16/slug_churn_m04q01_001_m07q20_001.png}
  \caption{Comparação pontual entre \textit{slug}, m04q01, e \textit{churn}, m07q20:
  regiões de coerência elevada também ocorreram entre padrões diferentes.}
  \label{fig:disc-inter}
  \fonte{Elaborado pelo autor, execução espacial 18--71.}
\end{figure}

A maximização por escala admitiu um atraso diferente para cada comprimento de onda
$\lambda$. O escore de coerência $S_{xy}^{(0{,}05)}$ não identificou, portanto,
um único deslocamento que alinhasse um mesmo evento nos dois registros.
Não se demonstrou uma causa única para a diferença de desempenho entre as medidas:
o domínio observado, a extensão dos prefixos e a forma de resumir os campos
também diferiram. Atribuir o resultado apenas à ausência de fase espacial comum
teria excedido o que foi verificado.

% =========================================================
\section{Sensibilidade ao piso de potência}\label{sec:disc-pisos}
% =========================================================

O piso de potência definiu quais escalas participaram do escore de coerência
$S_{xy}^{(\eta)}$, conforme a \autoref{eq:met-piso}. Ele não foi um limiar de
significância e não retirou regiões dos mapas. A \autoref{tab:disc-pisos}
reuniu os três pisos registrados na execução. Em todos, o escore de coerência
$S_{xy}^{(\eta)}$ médio foi menor dentro dos padrões do que entre padrões.
A direção do resultado, portanto, não se restringiu ao piso principal de 5\%.

\begin{table}[htbp]
  \centering
  \caption{Sensibilidade do escore de coerência $S_{xy}^{(\eta)}$ ao piso de potência,
  com peso igual por par.}
  \label{tab:disc-pisos}
  \begin{tabular}{c r r}
    \hline
    \textbf{Piso de potência $\eta$} &
    \shortstack{\textbf{Escore médio}\\\textbf{intrapadrão}} &
    \shortstack{\textbf{Escore médio}\\\textbf{entre padrões}}\\
    \hline
    1\% & 0,1907 & 0,2652\\
    5\% & 0,1877 & 0,2635\\
    10\% & 0,1822 & 0,2614\\
    \hline
  \end{tabular}
  \fonte{Elaborado pelo autor a partir do resumo intrapadrão e entre padrões da execução.}
\end{table}

Essa estabilidade de direção ajudou a delimitar o problema: alterar somente o piso
entre os valores examinados não produziu maior escore médio intrapadrão.
O resultado ofereceu uma referência para novas análises de resumo dos campos,
sem estabelecer que toda informação contida na coerência wavelet fosse inadequada
à caracterização dos padrões.

% =========================================================
\section{Qualidade da entrada e sensibilidade a m03q40}\label{sec:disc-qualidade}
% =========================================================

A auditoria reconstruiu os 31 panoramas de origem a partir de 15.531 quadros
e reproduziu suas intensidades pixel a pixel. Essa verificação confirmou a execução
do algoritmo de montagem. Ela não demonstrou que cada deslocamento estimado
representasse corretamente o movimento físico das interfaces.

No ensaio anular m03q40, mostrado na \autoref{fig:disc-m03q40}, 60 dos 500
deslocamentos estimados, ou 12\%, haviam sido substituídos por tiras de um pixel
porque eram nulos ou negativos. Além disso, 25,43\% dos pixels da faixa 18--71
do panorama atingiram a intensidade máxima de 255. O acúmulo no teto de intensidade
foi observado antes da padronização; sua origem não foi isolada entre aquisição,
processamento anterior e cena. A normalização posterior não recuperou informação
que eventualmente tivesse sido perdida nesse teto.

\begin{figure}[p]
  \centering
  \includegraphics[width=\textwidth,height=0.76\textheight,keepaspectratio]{figuras/v16/wavy_annular_m03q40_001.png}
  \caption{Panorama e sinais de m03q40, mantido na análise com ressalva de qualidade.
  A faixa clara do panorama acompanhou o acúmulo de pixels no teto de intensidade.}
  \label{fig:disc-m03q40}
  \fonte{Elaborado pelo autor, execução espacial 18--71.}
\end{figure}

O coeficiente de Pearson $r$ médio de m03q40 com os outros oito ensaios anulares
foi $0{,}7424$; entre esses outros oito, foi $0{,}9508$.
A \autoref{tab:disc-controles} mostrou os controles realizados para examinar essa
diferença. Em cada linha, a mesma regra foi aplicada aos nove ensaios.
O prefixo anterior à primeira substituição reduziu a discrepância, mas não
a eliminou. A faixa vertical mais estreita e a equalização também alteraram
os coeficientes de Pearson $r$, enquanto a retirada da remoção de tendência
linear produziu pouca mudança.

\begin{table}[htbp]
  \centering
  \caption{Coeficiente de Pearson $r$ médio nos controles do grupo anular.}
  \label{tab:disc-controles}
  \small
  \begin{tabular}{p{0.55\textwidth} r r}
    \hline
    \textbf{Controle} &
    \shortstack{\textbf{m03q40 versus}\\\textbf{outros oito}} &
    \shortstack{\textbf{Entre os}\\\textbf{outros oito}}\\
    \hline
    Referência: panoramas completos, linhas 18--71 & 0,7424 & 0,9508\\
    Prefixo comum de 14.008 px, anterior à primeira substituição & 0,8491 & 0,9350\\
    Comprimentos de onda de 50 a 1024 px & 0,8675 & 0,9465\\
    Sem remoção da tendência linear do sinal & 0,7405 & 0,9508\\
    Média das linhas 20--59 & 0,8754 & 0,9508\\
    Equalização global do histograma por aquisição & 0,8763 & 0,9249\\
    \hline
  \end{tabular}
  \fonte{Elaborado pelo autor a partir dos controles de m03q40 da execução.}
\end{table}

Os controles revelaram sensibilidade ao processamento e à região observada,
sem separar seus efeitos da variabilidade física do escoamento. O rótulo anular
concordou com a tabela manual, e m03q40 e m03q20 compartilharam as mesmas
configurações de câmera nos campos verificados. Não se identificou, portanto,
um ajuste exclusivo da câmera ou uma troca de rótulo que explicasse o desvio.
Tampouco se demonstrou que as substituições de deslocamento fossem sua única causa.

A \autoref{tab:disc-exclusao} apresentou a análise com e sem m03q40.
A retirada deixou 25 ensaios, 79 pares intrapadrão e 221 entre padrões.
O coeficiente de Pearson $r$ médio intrapadrão do anular aumentou de
$0{,}9045$ para $0{,}9508$, mas o contraste descritivo $\Delta_h$ global do coeficiente
de Pearson $r$ diminuiu de $0{,}0659$ para $0{,}0549$.
Isso ocorreu porque a média entre padrões também aumentou.
A maior homogeneidade do grupo anular não implicou maior contraste global.

\begin{table}[htbp]
  \centering
  \caption{Sensibilidade à retirada de m03q40, com peso igual por par.}
  \label{tab:disc-exclusao}
  \small
  \begin{tabular}{l r r r r}
    \hline
    \textbf{Medida} & \textbf{Ensaios} & \textbf{Intra} & \textbf{Entre} &
    \shortstack{\textbf{Contraste}\\\textbf{descritivo $\Delta_h$}}\\
    \hline
    Coeficiente de Pearson $r$ & 26 & 0,8697 & 0,8038 & 0,0659\\
    Coeficiente de Pearson $r$ & 25 & 0,8826 & 0,8277 & 0,0549\\
    Escore de coerência $S_{xy}^{(0{,}05)}$ & 26 & 0,1877 & 0,2635 & $-0{,}0759$\\
    Escore de coerência $S_{xy}^{(0{,}05)}$ & 25 & 0,1956 & 0,2687 & $-0{,}0730$\\
    \hline
  \end{tabular}
  \fonte{Elaborado pelo autor a partir da sensibilidade à exclusão de m03q40.}
\end{table}

O contraste descritivo $\Delta_h$ do coeficiente de Pearson $r$ permaneceu positivo,
enquanto o contraste descritivo $\Delta_h$ do escore de coerência
$S_{xy}^{(0{,}05)}$ permaneceu negativo. A direção dos achados principais
não dependeu, portanto, apenas desse registro. Manteve-se o conjunto de 26 ensaios
como referência, acompanhado da análise exploratória de exclusão, sem selecionar
aquisições pelo aumento da correlação.

% =========================================================
\section{Comparação com a literatura}\label{sec:disc-literatura}
% =========================================================

A \autoref{quad:disc-literatura} situou a contribuição executada em relação
a trabalhos que utilizaram transformadas wavelet ou classificação por imagens.
A comparação foi metodológica: não se equipararam frequências temporais em hertz
a comprimentos de onda do panorama em pixels, nem coeficientes de similaridade
a taxas de acerto de classificadores.

\begin{quadro}[htbp]
  \centering
  \caption{Relação entre os antecedentes bibliográficos e a análise executada.}
  \label{quad:disc-literatura}
  \small
  \begin{tabular}{p{0.22\textwidth} p{0.33\textwidth} p{0.35\textwidth}}
    \hline
    \textbf{Referência} & \textbf{Procedimento reportado} & \textbf{Relação com este trabalho}\\
    \hline
    \citeonline{nguyen2010} &
    Combinaram descritores da CWT, fração de vazio média e razão entre o comprimento
    efetivo estimado da estrutura gasosa e o diâmetro do tubo. &
    Fundamentaram o interesse em descritores espectrais por registro; aqui,
    foram examinados descritores de intensidade dos panoramas.\\
    \hline
    \citeonline{sun2012} &
    Combinaram entropia de norma de sinais de pressão com vazão mássica total
    ou fração volumétrica de gás em mapas de identificação. &
    Relacionaram descritores wavelet à identificação de padrões, com outra
    grandeza medida e outro descritor.\\
    \hline
    \citeonline{deoliveira2021} &
    Analisaram por CWT séries temporais do sensor óptico e de pressão diferencial,
    relacionando frequências à morfologia e às condições do escoamento. &
    Ofereceram um antecedente no arranjo experimental; a análise atual
    utilizou a coordenada espacial reconstruída.\\
    \hline
    \citeonline{dias2025knn} e \citeonline{dias2025cnn} &
    Classificaram imagens por atributos, incluindo DWT, com kNN e por rede
    convolucional ResNet50, respectivamente. &
    Avaliaram classificadores; aqui, quantificou-se a semelhança dos
    descritores, sem treinamento ou teste preditivo.\\
    \hline
  \end{tabular}
  \fonte{Elaborado pelo autor a partir das referências indicadas.}
\end{quadro}

Esses antecedentes sustentaram a investigação de representações espectrais e
visuais para os padrões. A presente execução acrescentou uma comparação
rastreável entre a forma dos resíduos espectrais e um resumo da coerência
alternativa, mantendo explícita a diferença entre caracterização descritiva
e reconhecimento automático.

% =========================================================
\section{Alcance dos achados e continuidade da análise}\label{sec:disc-limitacoes}
% =========================================================

As limitações da \autoref{quad:disc-limites} delimitaram a interpretação dos
resultados. Os grupos foram desbalanceados e cada aquisição participou de
vários pares. A média por grupo tornou a ponderação mais transparente,
mas não criou novas réplicas nem eliminou a dependência entre comparações.

\begin{quadro}[htbp]
  \centering
  \caption{Limitações observadas e consequências para a interpretação.}
  \label{quad:disc-limites}
  \small
  \begin{tabular}{p{0.40\textwidth} p{0.52\textwidth}}
    \hline
    \textbf{Condição observada} & \textbf{Consequência}\\
    \hline
    Três ensaios de \textit{plug}, três de \textit{slug} e um de bolhas dispersas &
    As médias dos grupos pequenos dependeram de poucas aquisições;
    não se estimou repetibilidade de bolhas dispersas.\\
    \hline
    Pares com aquisições compartilhadas e escalas correlacionadas &
    Nem os 325 pares nem as 200 escalas corresponderam a réplicas independentes.\\
    \hline
    Extensões de 3171 a 23079 px e faixa espectral fixa de 50 a 3000 px &
    Na maior escala, o menor panorama conteve apenas cerca de 1,06
    comprimentos de onda; o suporte variou entre registros.\\
    \hline
    Montagem por deslocamentos ópticos e ausência de calibração espacial &
    Os pixels descreveram a coordenada reconstruída; não se mediram
    diretamente dimensões físicas de bolhas ou velocidades de fase.\\
    \hline
    Exposições de 62, 127 e 132 $\mu$s e taxas de 4000, 5000 e 7000 quadros/s &
    A padronização de intensidade não eliminou possíveis diferenças
    de borramento, resposta óptica e teto de intensidade.\\
    \hline
    Condições operacionais distintas entre aquisições &
    Não se isolou experimentalmente o efeito do padrão dos efeitos
    das condições de operação e aquisição.\\
    \hline
  \end{tabular}
  \fonte{Elaborado pelo autor a partir dos dados e da auditoria da execução.}
\end{quadro}

Apesar dessas restrições, o contraste positivo do coeficiente de Pearson $r$
permaneceu nas ponderações examinadas e após a retirada exploratória de m03q40.
Esse resultado forneceu uma base favorável à continuidade da análise espectral,
com a extensão do panorama, a qualidade óptica e a composição dos grupos como
pontos concretos a controlar. A execução deixou essas questões abertas e
produziu dados rastreáveis para aprofundá-las, sem antecipar desempenho de
classificação ou significância estatística ainda não avaliados.
